
\documentclass[12pt]{article}
\usepackage[a4paper, margin=1in]{geometry}
\usepackage{parskip}
\usepackage{titlesec}
\titleformat{\section}{\large\bfseries}{\thesection}{1em}{}

\title{Presentation Script: Structural Loading on Appendages}
\date{}
\begin{document}

\maketitle

\section*{Slide 1: Title}
Welcome everyone. Today, we'll explore \textbf{Structural Loading on Appendages}, focusing on how wings and fins in nature manage aerodynamic forces. We'll discuss the physics behind it and walk through a practical example using the albatross takeoff scenario.

\section*{Slide 2: Motivation}
Why does this matter?
\begin{itemize}
    \item Nature offers elegant solutions for flight.
    \item Understanding these mechanisms helps us design efficient and safe aircraft.
    \item Structural loading is key to performance and avoiding failure.
\end{itemize}

\section*{Slide 3: Propulsive Appendages in Nature}
Let’s look at how animals use their appendages:
\begin{itemize}
    \item Fish use fins to steer and push through water.
    \item Birds and insects use wings to generate lift and control direction.
    \item These appendages must endure loads without breaking or deforming uncontrollably.
\end{itemize}

\section*{Slide 4: Gliding Ray and Lift Distribution}
This figure shows a gliding ray:
\begin{itemize}
    \item It uses its fins to glide, generating lift like a bird's wing.
    \item The figure on the right simplifies this concept using a twisted wing.
    \item This introduces the idea of lift distribution across a span.
\end{itemize}

\section*{Slide 5: 3D Effects – Root and Tip Vortices}
\begin{itemize}
    \item Real wings are not infinite.
    \item At the root and tip, 3D effects become important.
    \item Vortices form at the tips and alter lift distribution.
    \item Prandtl’s lifting-line theory helps model this behavior mathematically.
\end{itemize}

\section*{Slide 6: Bio-Inspired Propulsion}
This slide shows wings of a samara and a bio-inspired propeller blade:
\begin{itemize}
    \item Both are broken into 2D slices for analysis.
    \item The inflow velocity and rotation define how much force each slice feels.
\end{itemize}

\section*{Slide 7: Modeling Structural Forces}
We simplify by analyzing cross sections:
\begin{itemize}
    \item Forces include lift ($C_L$), drag ($C_D$), normal ($C_N$), and axial ($C_X$).
    \item Moments are bending in two directions and twisting (torsion).
\end{itemize}

\section*{Slide 8: Force Components on a Section}
This airfoil diagram helps explain the forces:
\begin{itemize}
    \item Lift acts upward, drag backward.
    \item From these we compute the normal and axial components that matter for structure.
\end{itemize}

\section*{Slide 9: Albatross Takeoff Scenario}
Let’s move into a real-world case:
\begin{itemize}
    \item The albatross takes off at a 30° angle with 15 m/s speed.
    \item Its wingspan and geometry are known.
    \item We analyze the forces and structural moments involved at the moment of lift-off.
\end{itemize}

\section*{Slide 10: Takeoff Modeling Example}
We model the wing using:
\begin{itemize}
    \item An elliptical planform
    \item Given angle of attack and velocity
    \item Using simple aerodynamic formulas, we compute:
    \begin{itemize}
        \item Lift and drag coefficients
        \item Resulting force components
        \item Flapwise, edgewise, and torsional moments
    \end{itemize}
\end{itemize}

\textbf{Conclusion:} These simplified models give us insight into how wings bear loads, showing the balance between aerodynamic efficiency and structural integrity.

\end{document}
